\documentclass[11pt]{article}
\usepackage{amsmath,amssymb}
\usepackage[hidelinks]{hyperref}
\title{Crossing a blocked relay region without breaking either link}
\date{}
\begin{document}
\maketitle

A relay has to move around a stopped robot while remaining within reach
of two fixed anchors. Feasible starting and ending positions do not
guarantee a feasible path. Even when a path exists in either of two
possible obstacle locations, the location has to be learned by a
deadline. This model gives an exact geometric threshold and an exact
latest-report criterion for a single speed-limited relay. It also
separates these results from a restricted acceleration comparison.

\section{Plant and obligations}
The idealised planar relay is a point, with obstacle clearance radius
\(\rho>0\) already inflated for the physical footprints. The anchors
are at \((\pm a,0)\), and both links require distance at most \(R>a>0\)
at all times. Assume this disk model; it does not cover shadowing by a
failed robot, packet scheduling, moving anchors, other obstacles or
steering and acceleration constraints unless separately specified.
Radio equality and boundary clearance count as lawful. The goal is to
move from \((0,-b)\) to \((0,b)\) while keeping both links and staying
outside the open obstacle disk, with speed at most \(v>0\).
A trajectory is absolutely continuous; stopping is allowed in this
kinematic model. Assume \(\rho\leq b\leq\sqrt{R^2-a^2}\), so both
endpoints can remain lawful. Abandoning one link temporarily is a
different service contract, not a recovery witness here.

\section{Centered obstruction and exact distance}

For a disk centered at zero, the relay lens has horizontal section
\([-(R-a),R-a]\) at \(y=0\). The intermediate value theorem forces any
crossing to visit that section. If \(\rho>R-a\), its entirety is
excluded. If \(\rho\leq R-a\), the whole obstacle circle lies in the
lens, and a shortest exterior path consisting of two tangent segments
and their joining circular arc stays in the lens by convexity.
Hence the exact connectivity-and-clearance threshold is
\(\rho\leq R-a\), and the exact minimum crossing length is
\[
       L_{\rm center}
       =2\sqrt{b^2-\rho^2}+2\rho\arcsin(\rho/b).
\]
The deadline threshold is \(L_{\rm center}/v\). At equality of
clearance and lens half width, the path uses a point with no slack
in a consumed physical bound. The equality case is not a robust
field guarantee.

Here is a length argument also used below. Measure polar radius and
angle from the obstacle center. Let \(m\) be the minimum radius on an
admissible trajectory and \(d_s,d_g\) its endpoint radii. If the minimum
angular separation of the endpoints is \(\Delta\), the total angular
variation is at least \(\Delta\). Pointwise Cauchy--Schwarz bounds the
speed from below by
\[
       \sqrt{1-(m/r)^2}\,|\dot r|+m|\dot\theta|.
\]
Integrating from each endpoint radius down to \(m\) bounds the total
length from below by
\[
 F(m)=\sum_{d\in\{d_s,d_g\}}
       \left(\sqrt{d^2-m^2}-m\arccos(m/d)\right)+m\Delta ,
\]
with the two summands counted separately if the endpoint radii agree.
The polar lift is absolutely continuous after arc-length
parametrization because the radius stays positive. For the centered
case \(\Delta=\pi\) and \(F\) increases with \(m\); tangent segments
and the shorter connecting arc attain \(F(\rho)\).
More generally, for endpoints \((0,-h)\), \((0,b)\) and center
\((c,0)\), take \(0<c<\rho\) and \(h,b\geq\sqrt{\rho^2-c^2}\).
Then \(d_s=\sqrt{c^2+h^2}\), \(d_g=\sqrt{c^2+b^2}\), and
\(\Delta=\arctan(h/c)+\arctan(b/c)\).
Here \(F'(m)=\Delta-\arccos(m/d_s)-\arccos(m/d_g)>0\)
for \(m\geq\rho\), using \(m>c\).
For a known center with \(0<c<\rho\) and \(c\leq a\),
a crossing exists exactly when \(\rho-c\leq R-a\), given
legal endpoints. Necessity uses the horizontal section as
before: the nearest allowed side of the obstacle starts
at \(x=c-\rho\), to the left of its center. For sufficiency,
the minimizing arc lies in \(x\leq0\), not just in the far
semicircle. For an endpoint at height magnitude \(h\),
the angular deviation of its tangency from the negative
horizontal ray is
\[
     \beta(h)=\arctan(h/c)-\arccos\bigl(\rho/\sqrt{c^2+h^2}\bigr).
\]
It is nonnegative and decreases for \(h>\sqrt{\rho^2-c^2}\)
by differentiation. At equality its value is
\(\arccos(c/\rho)\), which places the rightmost relevant
arc point at \(x=0\). On the arc the distance to the right
anchor dominates the left one, and is largest at the
left tip, at distance \(a-c+\rho\). Thus the arc fits at the
necessary radius and the tangent segments fit by convexity.
Define \(L(h)=F(\rho)\) for these endpoints; this is the exact
known-center suffix distance, not a sampled-path approximation.
The entire far semicircle need not fit.

\section{An exact deadline to learn the side}

Now the physical obstacle center is either \((c,0)\) or \((-c,0)\).
Exactly one exists. Keep the known-world conditions above and assume
\(\rho+c>R-a\). The relay receives a trustworthy sign at a fixed
time \(\tau\); before that, all its admissible observations are the
same in the two worlds. No probing action that reveals the sign
before \(\tau\) is part of this model. Let
\[
 h_*=\sqrt{\rho^2-c^2},\quad
 h(t)=\max(h_*,\,b-vt),\qquad
 A(t)=t+L(h(t))/v .
\]
The minimum worst-world arrival time among all causal trajectories
under this law is exactly \(A(\tau)\). Thus a guaranteed deadline
\(T\) is attainable precisely when \(A(\tau)\le T\).

For the construction, move vertically to \((0,-h_*)\) as far as
time allows, wait if necessary, and after the report take the
appropriate tangent-and-arc path from \((0,-h(\tau))\).
The common vertical segment is outside both possible open disks
and in the lens. The known-world path is admissible by the
far-arc bound and symmetry.

For the lower bound, every pre-report trajectory must avoid
both possible disks. It cannot yet cross \(y=0\), since no point
of that line in the lens is allowed in both worlds. If its
position at report time has \(x\geq0\), consider the actual
right-hand disk at \(c\). A crossing of \(y=0\) in this world
must have \(x<0\). Therefore the post-report path reaches
\(x=0\) at some time \(t'\geq\tau\) while still below the axis.
If it starts this part already on \(x=0\), take \(t'=\tau\).
That crossing height has depth at least \(h_*\) by clearance
and at least \(b-vt'\) by the total speed bound. Its remaining
distance is at least \(L(h(t'))\). To see that \(L(h)\) is
increasing, differentiate its closed form: the derivative
contributed by the lower endpoint is
\((h\sqrt{c^2+h^2-\rho^2}+\rho c)/(c^2+h^2)>0\).
Also \(A(t)\) is nondecreasing by the triangle inequality along
the safe central line: between two staging depths the known-world
distance can improve by at most the vertical distance moved.
Consequently arrival in this world is at least \(A(t')\geq
A(\tau)\). The reflected argument covers \(x<0\). Waiting
off center does not improve the worst case. If no sign arrives,
the common section obstruction persists independently of deadline.
This is an information obstruction, not physical impossibility
in either known world.

For example, with \(a=3,R=5,c=0.5,\rho=1.75,b=3.5,v=1\),
each known world takes about 7.4598 time units from the start.
For a deadline of 9 the sign must arrive by about 3.0283
time units. These are evaluations of exact expressions, not
measured flight or radio times.

\section{A separate bounded-curvature comparison}
Under a centered disk, consider only constant-speed tangent--arc
paths with turning radius \(r\in[\rho,\min(b,R-a)]\), speed
limit \(v\) and acceleration magnitude at most \(A_{\max}>0\)
almost everywhere. End velocities are allowed to be the
corresponding tangents; this experiment does not require rest
or bound jerk. Along a circular arc, speed may be at most
\(\min(v,\sqrt{A_{\max}r})\). The best time in this family is
the centered length at radius \(r\) divided by that speed.
In the acceleration-limited part, the derivative changes sign
exactly when \(u\tan u=1\), where \(r=b\sin u\).
In the speed-limited part the length increases with \(r\).
Thus the optimum radius is the clipping of
\(\min(b\sin u_*,v^2/A_{\max})\) to the allowed interval,
where \(u_*\) is the unique root in \((0,\pi/2)\).
A wider turn can be faster under this control law.
It is not an optimum over variable-speed acceleration-controlled
trajectories or a model of the hidden-case stop and branch dynamics.

\section{Checks and comparison}
\href{https://doi.org/10.1016/j.jfranklin.2023.05.013}{Maeda,
Endo and Matsuno} study decentralised connectivity and recovery
after an individual failure. This is a separate exact reachability
benchmark, not an implementation or performance test of their
controller. \href{https://arxiv.org/abs/2609.29939}{Sekercioglu,
De Carli and Dimarogonas} allow a formation manager to establish
prospective edges and temporarily relax their distance bounds.
Such a relaxation is not an admissible action under the two
fixed links here. \href{https://arxiv.org/abs/2210.01354}{Breeden
and Panagou} address feasibility of jointly applied barrier
conditions; a feasible command or invariant component alone
does not prove this target is reachable in time.
\href{https://epubs.siam.org/doi/10.1137/21M1428248}{Hershberger,
Suri and Yildiz} give the general curved-obstacle shortest-path
context. At this exact plant and information, independently
integrating the radial bound and a true-arc visibility-graph
reference agree in the tested known worlds. The graph
approximation checks examples, not the continuous obstruction.
Lean checks the algebraic crossing and containment kernels,
not polar path regularity, source validity or the full control
theorem. This is a scoped R6 synthesis, not I2, field assurance
or a claim of improvement over a matched exact planner.
\end{document}
